% LSPV-4 REPORT::2048F::2026-09-30T20:48:00-07:00
\documentclass[11pt]{article}
\usepackage{amsmath,amssymb}
\begin{document}
\section*{LSPV-4 red-team report 2048F}
\textbf{STATUS: OPEN.} P versus NP remains open.

\paragraph{PROVED/EXACT: proof-length feasibility barrier for accepted-family counting.}
Let $L$ be the PCP proof length and
\[
\Psi:\mathcal P_{\rm acc}\to\mathcal F_{\rm acc}
\]
map perfectly accepted proof strings to distinct extensional oracle objects at a fixed RM node. Necessarily,
\[
|\mathcal F_{\rm acc}|\le |\mathcal P_{\rm acc}|\le 2^L.
\]
LSPV-2's sufficient counting criterion for forcing maximal quadratic-feature rank has the form
\[
|\mathcal F_{\rm acc}|>2^{\kappa+M},
\qquad
M=1+m+\binom m2.
\]
Therefore that strategy is feasible only if
\[
\boxed{L>\kappa+M.}
\]

For the checked Theorem-24 parameterization, the proof length is
\[
L=T(n)\operatorname{polylog}T(n),
\]
the construction chooses $h=m^{100}$, $r$ as the least power of two with $h^r>L$, and $\Delta=hr$, while the local-extractor seed parameter satisfies
\[
\kappa=\Delta^5=(m^{100}r)^5=m^{500}r^5.
\]
Hence accepted-family counting cannot cross its own threshold whenever
\[
\boxed{
T(n)\operatorname{polylog}T(n)
\le
m^{500}r^5+1+m+\binom m2.
}
\]

This does \emph{not} falsify the counting theorem. It gives a necessary feasibility condition that must be checked before using that theorem as an active attack on quadratic contraction. Collisions in $\Psi$ only strengthen the obstruction.

\paragraph{SAME\_REALIZATION interpretation.}
In the checked vMS construction, a successful computation path guesses the target value and proof data, verifies perfect acceptance, and then emits the RMV object attached to that accepted realization. Thus accepted proofs legitimately parameterize successful realizations, but there can be no more distinct extensional node oracles than accepted proof strings. Proof padding, COPY, source relabeling, or duplicate Merlin strategies do not evade the upper bound $|\mathcal F_{\rm acc}|\le 2^L$.

\paragraph{Weakest repair.}
Define
\[
\Theta(n,m)
=
L-(\kappa+M)
=
T(n)\operatorname{polylog}T(n)
-
\left(
m^{500}r^5+1+m+\binom m2
\right).
\]
Before any accepted-family cardinality argument, compute the sign of $\Theta$ in the contradiction regime.

If $\Theta\le0$, retire accepted-family counting for that parameter regime and return to direct Schur-square/product-span or explicit independent-minor analysis on a single legal realization.

If $\Theta>0$, the remaining obligation is sharper:
\[
\log_2|\operatorname{im}\Psi|>\kappa+M
\]
must be proved for \emph{distinct extensional node oracles}, not merely for accepting proof strings.

\paragraph{Relation to LSPV-1 controls.}
LSPV-1's finite nonconstant-leaf full-rank experiments remain \textbf{COMPUTATIONAL} and logically independent of this accepted-family threshold. They do not rely on counting accepted proofs. Their all-parameter lift and all-zero accepted-proof control remain \textbf{OPEN}.

\paragraph{Other active attacks.}
Exact canonical $O(m)$ product span remains \textbf{OPEN}. An architecture-preserving $\Omega(m^2)$ minor remains \textbf{OPEN}. Sparse-$Q$ acquisition and coefficient-compatible size normalization remain \textbf{OPEN}. Quantitative $k^3$ SAT-prefix absorption remains \textbf{OPEN}, while the conservative $k^4$ reconstruction ledger remains \textbf{CONDITIONAL}. TransportCert/atomic TC4 and ordinary switching remain \textbf{FALSIFIED/RETIRED} as universal lower-bound bridges. SWITCH+SURVIVE+BIND remains a \textbf{CANDIDATE} only if BIND becomes nonextensional, wrapper-invariant, and unavoidable under maximal semantic and algebraic oracles.

\paragraph{Barrier audit and P=NP evidence.}
The counting barrier itself relativizes and algebrizes. A maximal semantic oracle can reveal the accepted image, its size, rank, or a sparse annihilator. This therefore supplies no missing nonrelativizing or nonalgebrizing hinge. There is \textbf{no genuine positive evidence for $P=NP$} in this run. Failure of a proposed $P\ne NP$ mechanism is not evidence for equality.

\paragraph{Strict status.}
PROVED/EXACT: accepted-family proof-length feasibility barrier $L>\kappa+M$.
CONDITIONAL: the LSPV-2 counting theorem as a route to maximal rank, subject to image-size feasibility.
COMPUTATIONAL: LSPV-1 finite full-rank controls.
OPEN: exact canonical product span; canonical $\Omega(m^2)$ minor; sparse-$Q$ acquisition and size normalization; $k^3$ flattening; wrapper-invariant BIND; relativization/algebrization escape; final diagonal.
FALSIFIED/RETIRED: ordinary switching and unrestricted TC4 as universal lower-bound bridges.

\paragraph{Single next falsification target.}
Substitute the actual contradiction regime's $T(n)$ and target circuit size $m(n)$ into
\[
\Theta(n,m)
=
T(n)\operatorname{polylog}T(n)
-
\left(
m^{500}r^5+1+m+\binom m2
\right).
\]
If $\Theta\le0$ eventually, LSPV-2's accepted-family cardinality strategy is rigorously unavailable in that regime. If $\Theta>0$, the next red-team target is the exact image size of the accepted-proof-to-node-oracle map $\Psi$, not proof multiplicity itself.

\begin{thebibliography}{3}
\bibitem{SU05} R. Shaltiel and C. Umans, J. ACM 52(2), 2005.
\bibitem{vMS25} D. van Melkebeek and N. Mocelin Sdroievski, Computational Complexity 34:15, 2025.
\bibitem{AW09} S. Aaronson and A. Wigderson, ACM TOCT 1(1), 2009.
\end{thebibliography}
\end{document}
% END_OF_VERIFIED_REPORT::LSPV-4::2048F